\chapter*{Výchozí úvaha autora}
\addcontentsline{toc}{chapter}{Výchozí úvaha autora}

Mým cílem je zaznamenávat pohyb formule ve 3D prostoru, především přesně
vůči základně; přibližné umístění na mapě mi zpočátku postačuje.
Chci zachovat surová data GNSS, surová data IMU a průběžný kombinovaný
výstup pro porovnání online řešení s postprocessingem.

Pro GNSS/IMU část si stanovuji přibližně 25~000~Kč bez DPH s možností
malého překročení. SSD dodám vlastní, notebook využiji při vývoji.
Cílově chci základnu ZED-F9P s ESP32, raw loggerem a webovou správou
bez notebooku. Rychlý mosaic-X5 zachovávám na formuli.

Nechci pořizovat drahou taktickou IMU bez opory v měření. Pro první návrh
volím rozsah \SI{\pm16}{g} a vzorkování kolem \SI{500}{Hz}. Informaci o
nárazové ochraně kolem \SI{3}{g} beru jako vodítko, nikoli záruku proti
lokálním špičkám. Rozestup vzorků nezaměňuji s přesností polohy.

Původní ISM330DHCX nahrazuji deskou Adafruit ICM-20948 kvůli kompasu
a dostupnosti u RPiShopu pro společnou objednávku s Pi a rádiovými sadami.
Musím ověřit šum, drift a magnetické rušení. Omezení životnosti čipu
uvádím v kapitole o výběru.

IMU chci připojit přímo k Raspberry Pi 5 s Raspberry Pi OS, bez časovacího
STM32 či Teensy. Přesnost časového přiřazení musím změřit. ESP32 v rádiových
sadách bude obsluhovat korekční spoj, nikoli sběr IMU.

Pro korekce chci přímé LoRa spojení 868~MHz: dvě sady pro spoj, dvě
pro vývoj a rezervu. Ověřím propustnost, rádiové podmínky a pokrytí.
Dvě ANN-MB2 volím pro vlastní kryty a úsporu; montáž ověřím měřením.
NTRIP ponechávám jako zálohu.

Po konzultaci se školitelem upřednostňuji dostupnost a méně dodavatelů před
úsporou jednotlivých stokorun. Volím RPiShop a ArduSimple; Mouser je
nákupní alternativa. Součet přibližně 24~290~Kč bez DPH nezahrnuje
napájení, úložiště základny, instalační díly ani dopravu. Vlastní přínos
vidím v integraci, firmwaru a ověření, nikoli v nutnosti vyvíjet GNSS čip.

V navazující etapě chci tento systém začlenit do senzoriky celé formule.
Raspberry Pi má ukládat i pohyb roll/heave, otáčky kol a teploty.
Nově volím centrální ADC u Raspberry a samostatné přívody od snímačů.
CAN 1 a CAN 2 ponechávám pro auto a ADU; místní uzly zatím nepořizuji.
Tento širší rozsah i jeho napájení posuzuji
odděleně od původního rozpočtu; nejde zatím o realizovanou instalaci.

\begin{flushright}
  \textit{Jiří Buršík, autor práce}
\end{flushright}
